\documentclass[12pt,twoside]{report}

%%%%%%%%%%%%%%%%%%%%%%%%%%%%%%%%%%%%%%%%%%%%%%%%%%%%%%%%%%%%%%%%%%%%%%%%%%%%%

% Definitions for the title page
% Edit these to provide the correct information
% e.g. \newcommand{\reportauthor}{Timothy Kimber}

\newcommand{\reporttitle}{Analysing property sales data using Data Science}
\newcommand{\reportauthor}{Wenxiang Luo}
\newcommand{\supervisor}{Chiraag Lala}
\newcommand{\degreetype}{MSc Computing}

%%%%%%%%%%%%%%%%%%%%%%%%%%%%%%%%%%%%%%%%%%%%%%%%%%%%%%%%%%%%%%%%%%%%%%%%%%%%%

% load some definitions and default packages
%%%%%%%%%%%%%%%%%%%%%%%%%%%%%%%%%%%%%%%%%
% University Assignment Title Page 
% LaTeX Template
% Version 1.0 (27/12/12)
%
% This template has been downloaded from:
% http://www.LaTeXTemplates.com
%
% Original author:
% WikiBooks (http://en.wikibooks.org/wiki/LaTeX/Title_Creation)
%
% License:
% CC BY-NC-SA 3.0 (http://creativecommons.org/licenses/by-nc-sa/3.0/)
% 
%
%%%%%%%%%%%%%%%%%%%%%%%%%%%%%%%%%%%%%%%%%
%----------------------------------------------------------------------------------------
%	PACKAGES AND OTHER DOCUMENT CONFIGURATIONS
%----------------------------------------------------------------------------------------
\usepackage[a4paper,hmargin=2.8cm,vmargin=2.0cm,includeheadfoot]{geometry}
\usepackage{textpos}
\usepackage[numbers]{natbib} % for bibliography
\usepackage{tabularx,longtable,multirow,subfigure,caption}%hangcaption
\usepackage{fncylab} %formatting of labels
\usepackage{fancyhdr} % page layout
\usepackage{url} % URLs
\usepackage[english]{babel}
\usepackage{amsmath}
\usepackage{graphicx}
\usepackage{dsfont}
\usepackage{epstopdf} % automatically replace .eps with .pdf in graphics
\usepackage{backref} % needed for citations
\usepackage{array}
\usepackage{latexsym}
\usepackage[pdftex,pagebackref,hypertexnames=false,colorlinks]{hyperref} % provide links in pdf

\hypersetup{pdftitle={},
  pdfsubject={}, 
  pdfauthor={},
  pdfkeywords={}, 
  pdfstartview=FitH,
  pdfpagemode={UseOutlines},% None, FullScreen, UseOutlines
  bookmarksnumbered=true, bookmarksopen=true, colorlinks,
    citecolor=black,%
    filecolor=black,%
    linkcolor=blue,%
    urlcolor=black}

\usepackage[all]{hypcap}


%\usepackage{color}
%\usepackage[tight,ugly]{units}
%\usepackage{float}
%\usepackage{tcolorbox}
%\usepackage[colorinlistoftodos]{todonotes}
% \usepackage{ntheorem}
% \theoremstyle{break}
% \newtheorem{lemma}{Lemma}
% \newtheorem{theorem}{Theorem}
% \newtheorem{remark}{Remark}
% \newtheorem{definition}{Definition}
% \newtheorem{proof}{Proof}


%%% Default fonts
\renewcommand*{\rmdefault}{bch}
\renewcommand*{\ttdefault}{cmtt}



%%% Default settings (page layout)
\setlength{\parindent}{0em}  % indentation of paragraph

\setlength{\headheight}{14.5pt}
\pagestyle{fancy}
\renewcommand{\chaptermark}[1]{\markboth{\chaptername\ \thechapter.\ #1}{}} 

\fancyfoot[ER,OL]{\sffamily\textbf{\thepage}}%Page no. in the left on odd pages and on right on even pages
\fancyfoot[OC,EC]{\sffamily }
\renewcommand{\headrulewidth}{0.1pt}
\renewcommand{\footrulewidth}{0.1pt}
\captionsetup{margin=10pt,font=small,labelfont=bf}


%--- chapter heading

\def\@makechapterhead#1{%
  \vspace*{10\p@}%
  {\parindent \z@ \raggedright \sffamily
    \interlinepenalty\@M
    \Huge\bfseries \thechapter \space\space #1\par\nobreak
    \vskip 30\p@
  }}

%---chapter heading for \chapter*  
\def\@makeschapterhead#1{%
  \vspace*{10\p@}%
  {\parindent \z@ \raggedright
    \sffamily
    \interlinepenalty\@M
    \Huge \bfseries  #1\par\nobreak
    \vskip 30\p@
  }}

\allowdisplaybreaks

% load some macros
% Here, you can define your own macros. Some examples are given below.

\newcommand{\R}[0]{\mathds{R}} % real numbers
\newcommand{\Z}[0]{\mathds{Z}} % integers
\newcommand{\N}[0]{\mathds{N}} % natural numbers
\newcommand{\C}[0]{\mathds{C}} % complex numbers
\renewcommand{\vec}[1]{{\boldsymbol{{#1}}}} % vector
\newcommand{\mat}[1]{{\boldsymbol{{#1}}}} % matrix


\date{June 2022}

\begin{document}

% load title page
% Last modification: 2015-08-17 (Marc Deisenroth)
\begin{titlepage}

\newcommand{\HRule}{\rule{\linewidth}{0.5mm}} % Defines a new command for the horizontal lines, change thickness here


%----------------------------------------------------------------------------------------
%	LOGO SECTION
%----------------------------------------------------------------------------------------

\includegraphics[width = 4cm]{./figures/imperial}\\[0.5cm] 

\center % Center remainder of the page

%----------------------------------------------------------------------------------------
%	HEADING SECTIONS
%----------------------------------------------------------------------------------------

\textsc{\Large Imperial College London}\\[0.5cm] 
\textsc{\large Department of Computing}\\[0.5cm] 

%----------------------------------------------------------------------------------------
%	TITLE SECTION
%----------------------------------------------------------------------------------------

\HRule \\[0.4cm]
{ \huge \bfseries \reporttitle}\\ % Title of your document
\HRule \\[1.5cm]
 
%----------------------------------------------------------------------------------------
%	AUTHOR SECTION
%----------------------------------------------------------------------------------------

\begin{minipage}{0.4\textwidth}
\begin{flushleft} \large
\emph{Author:}\\
\reportauthor % Your name
\end{flushleft}
\end{minipage}
~
\begin{minipage}{0.4\textwidth}
\begin{flushright} \large
\emph{Supervisor:} \\
\supervisor % Supervisor's Name
\end{flushright}
\end{minipage}\\[4cm]


%----------------------------------------------------------------------------------------
%	FOOTER & DATE SECTION
%----------------------------------------------------------------------------------------
\vfill % Fill the rest of the page with whitespace
Submitted in partial fulfillment of the requirements for the MSc degree in
\degreetype~of Imperial College London\\[0.5cm]

\makeatletter
\@date 
\makeatother


\end{titlepage}



% page numbering etc.
\pagenumbering{roman}
\clearpage{\pagestyle{empty}\cleardoublepage}
\setcounter{page}{1}
\pagestyle{plain}

%%%%%%%%%%%%%%%%%%%%%%%%%%%%%%%%%%%%
%\begin{abstract}
%Your abstract.
%\end{abstract}

\cleardoublepage
%%%%%%%%%%%%%%%%%%%%%%%%%%%%%%%%%%%%
%\section*{Acknowledgments}
%Comment this out if not needed.

\clearpage{\pagestyle{empty}\cleardoublepage}

%%%%%%%%%%%%%%%%%%%%%%%%%%%%%%%%%%%%
%--- table of contents
\fancyhead[RE,LO]{\sffamily {Table of Contents}}
\tableofcontents 


\clearpage{\pagestyle{empty}\cleardoublepage}
\pagenumbering{arabic}
\setcounter{page}{1}
\fancyhead[LE,RO]{\slshape \rightmark}
\fancyhead[LO,RE]{\slshape \leftmark}

%%%%%%%%%%%%%%%%%%%%%%%%%%%%%%%%%%%%
\chapter{Introduction}

Nowadays, there is a substantial amount of data generated every second. The daily lives of humans are producing it, and some other fields, such as research, health care, economic activities, and environmental information from various sensors, also generate a vast amount of data. Obtaining the relationship between some features or the patterns underlying these massive amounts of data might benefit the entire world. For instance, new causes of diseases might be identified, and technological advancement could be accelerated.
\\

However, this extensive data can be one of the main obstacles for analysis as it is approximately impossible for humans to obtain insights into the data manually. Under this circumstance, artificial intelligence (AI), a technique that empowers the computer to imitate human intelligence and manner,  could be one of the methods to mitigate this issue. It can extract patterns from large datasets and use them to make predictions based on future data and even identify which data components are responsible for the results.
\\

In this project, some AI techniques will be applied to property and demographic data to gain insights and understand the factors influencing a homeowner's likelihood to sell. The factors might include the proximity to schools, hospitals, or supermarkets, the accessibility to public transportation, and the property types (flats or houses). It could be highly advantageous to estate agents who would discover homeowners with more potential to become clients and provide them with business. 

%%%%%%%%%%%%%%%%%%%%%%%%%%%%%%%%%%%%
\chapter{Literal Review}
Python is one of the most popular programming languages in the world since it is simple to develop, and there are extensive packages for various functionalities. In this project, Python and its packages would be used for loading data, preprocessing data, and constructing and evaluating machine learning models. 

\section{Machine Learning}
Machine Learning (ML), a subset of AI, is a technique that the computer can learn and improve from data without explicit programming. The reason for utilizing ML is that its performance is sometimes better than the conventional approach. For example, ML techniques would simplify the solution to a problem that comprises a long list of rules (spam mail detection). 
\\

ML can be divided into three categories, one of which is supervised learning. In supervised learning, the dataset contains features (input to the model) and targets (ground truth of the output), and the model's parameters are randomly initialized. Then the features are passed to the model, and the differences between the current output and the ground truth are used to update the parameters until the differences are acceptable. 
\\

In this project, a supervised learning model will be implemented for data analysis, and the steps are listed below.
\begin{enumerate}
	\item Data Preprocessing: Some data from the dataset may be missing, and these values must be handled appropriately before being passed to the model. 
	\item Standardization: In real life, different features usually have different ranges, and this will cause a problem in ML, which is that high magnitude features would have more weight than low magnitude features \citep{RN4}. One of the solutions is standardization, which could scale all the features to the same magnitude.
	\item Feature encoding: ML models require numerical values, whereas the categorical features in the dataset do not satisfy this requirement. Therefore, these features should be converted into numerical values. 
	\item Training \& Testing: The parameters of the model are updated, and it is expected that the loss will converge during training. The performance of the model is validated when testing.
\end{enumerate}

\section{NumPy}
Numerical Python (\textit{\textbf{NumPy}}) is a scientific computing package utilizing an optimized C/C++ API to reduce computation time compared to pure Python computations \citep{RN6}. It provides some data structures, one of the most important structures is \textbf{\emph{ndarray}}. Unlike Python lists, NumPy arrays are homogenous, meaning all the elements in them are of the same type, and their sizes are constant \citep{RN4}.  In addition, NumPy provides various built-in functions for a variety of purposes, including statistics, linear algebra,  transforms, and element-wise operation. 
\\

In this project, \textit{\textbf{Numpy}} will be used for data preprocessing. For example, converting categorical data into integers.

\section{Pandas}
The \textbf{\textit{pandas}} is a fast and compelling package capable of handling data of various types (numerical values, strings, and time) and from multiple sources (CSV, Excel, and MqSQL database). One of the \textbf{\textit{pandas}} data structures is \textbf{\textit{DataFrame}} which is appropriate for handling tabular data with columns of different types. In addition, it could manage various operations, such as manipulating missing values, creating pivot tables (table summaries), and grouping data from different columns \citep{RN4}. 
\\

\textbf{\textit{Pandas}} would be used for data loading and preprocessing in this project. Data is initially loaded and represented as a \textbf{\textit{DataFrame}}. Then the attributes of the data, such as distribution, should be inspected. Finally, the data is preprocessed.

\section{Scikit-learn}
Scikit-learn is a Python ML package that provides various techniques for data mining, modeling, and analyzing. It could be commonly used in multiple ML tasks, such as classification, regression, and clustering. This project will use the package to analyze the performance of the models.

\section{PyTorch}
\textbf{\textit{PyTorch}} was developed by Facebook, and it is one of the popular libraries for constructing ML models. It is also the basis of numerous packages that the developers could take advantage of \citep{RN5}. In addition, this library is capable of performing auto differentiation by using graphic processing unit (GPU) acceleration, resulting in a reduction in training time. 
\\

In this project, some of the \textbf{\textit{PyTorch}} components will be utilized.
\begin{itemize}
	\item \textbf{\textit{nn.Sequential}} is the container that contains all the layers of the ML model
	\item An appropriate loss function would be selected to compute the loss and perform backpropagation. 
	\item One of the optimizers would be used to update the model's parameters. 
\end{itemize}

\section{Matplotlib}
\textbf{\textit{Matplotlib}} is one of the Python plotting packages for plotting different types of figures, including histograms, pie plots, and line plots. This library will be utilized primarily for data visualization in this project. For instance, the distribution of the raw data should be visualized to determine the procedures of preprocessing and plotting the training loss and validation loss during model evaluation. 



%%%%%%%%%%%%%%%%%%%%%%%%%%%%%%%%%%%%
\chapter{Project Plan}

\section{Data Preprocessing}
\subsection{Handling Missing Values}
In this project, two methods will be attempted to preprocess the missing values.
\begin{enumerate}
	\item Drop missing values: This method deletes the messing values for data analysis.
	\item Fill the missing values: In this approach, the missing values are filled with mean, median, or default values, for example, zero or some constants.
\end{enumerate}

\subsection{Standardization}
In this project, two different approaches will be attempted. 
\begin{enumerate}
	\item Standard scaling: By applying this method, the mean of the feature is removed and then divided by the standard deviation. 
	\item Min-max scaling: For this approach, the minimum and maximum of the raw data are used to transform it into a specific range.
\end{enumerate}

\subsection{Feature Encoding}
In this project, label encoding will be applied for this goal. It would convert the categorical values into a sequence of integer values. 

\subsection{Separating Dataset}
The original dataset should be shuffled and divided into three parts which will be used for training, validation, and testing. 


\section{Build Model with PyTorch}
\subsection{Model architecture}
\begin{itemize}
	\item \textbf{Basic Model}: The number of the input neural of the model should be the same as the number of features, and then there are hidden layers. Finally, the output layer is a sigmoid function as this model should output the possibility of the selling. 
	\item \textbf{Advanced Model}: The advanced model contains more components, such as dropout or batch normalization, and these additional layers might reduce overfitting and enhance accuracy. 
\end{itemize}

\subsection{Training}
\begin{enumerate}
	\item Use \textbf{\textit{DataLoader}} to fetch batches of data for training, and the batch  size could be set to 128 if there is sufficient memory. 
	\item The loss function used for model training is \textbf{\textit{MSELoss}} since this is a regression problem. 
	\item The optimizer for updating parameters could be \textbf{\textit{SGD}} or \textbf{\textit{Adam}}.
	\item Train the model iteratively, and the number of iterations (epochs) will be set to 500 and updated if necessary. 
\end{enumerate}

\subsection{Hyperparameters Tuning}
\begin{itemize}
	\item Learning rate: The initial learning rate will be set to 0.001, and it will be increased/decreased depending on the performance. Moreover, learning rate schedulers, for example, \textbf{\textit{LambdaLR}} and \textbf{\textit{StepLR}}, could be applied during training. 
	\item The number of epochs: If the initial value causes overfitting or underfitting, then the epochs should be decreased or increased. 
	\item Activation functions: The initial activation function will be ReLU, but other activation functions, such as TanH and Parametric ReLU, will also be tested to enhance the performance. 
\end{itemize}

\section{Evaluation}
After constructing and training models, the next step is to evaluate their performance. The mean squared error (MSE) could be used to evaluate performance numerically. The lower the MSE, the higher the performance. In addition, a line plot containing the training loss and validation loss against epochs should be produced. This figure could determine if the model is overfitting/underfitting. The optimal best model is then selected and tested using the test dataset to generate the final result. 


%%%%%%%%%%%%%%%%%%%%%%%%%%%%%%%%%%%%
%\chapter{Experimental Results}


%%%%%%%%%%%%%%%%%%%%%%%%%%%%%%%%%%%%
%\chapter{Conclusion}


%% bibliography
\bibliographystyle{apalike}
\bibliography{proposal}

\end{document}
